\chapter{Weakness}

\section*{Definition}
Decreased muscle power strength perform voluntary movements distinguish fatigue subjective tiredness objective reduced force generation measured medical examination physical therapy assessment

\section*{What Happens in Your Body}
Differentiate proximal distal symmetric asymmetric pattern determines differential diagnosis neuromuscular endplate junction disorders peripheral nerve root spinal cord neuromuscular junction muscular metabolism metabolic derangement electrolyte imbalance central nervous system lesions neuromuscular transmission defects

\section*{Causes}
\begin{itemize}
  \item Neuromuscular myasthenia gravis Lambert-Eaton syndrome
  \item Myopathies inflammatory polymyositis inclusion body genetic muscular dystrophy
  \item Peripheral neuropathies diabetic alcohol compressive radiculopathies spinal stenosis
  \item Central processes stroke tumor neurodegeneration ALS Parkinson multiple sclerosis
  \item Metabolic electrolyte abnormalities hypokalemia hyperkalemia hypocalcemia thyroid dysfunction anemia
\end{itemize}

\section*{When to See a Doctor}
See your doctor if:
\begin{itemize}
  \item Symptoms are severe or getting worse over time
  \item Weakness is progressive, unexplained, or affects your ability to walk or breathe
  \item You have other concerning symptoms such as fever, weight loss, or bleeding
\end{itemize}

\section*{Related Symptoms}
\begin{itemize}
  \item Fatigue exertion pattern muscle atrophy fasciculations reflexes diminished gait clumsiness balance problems falls
\end{itemize}
\textbf{Important:} If this symptom persists or concerns you, your doctor will check for serious weakness

\section*{Self-Care / Home Management}
\begin{itemize}
  \item Prioritize adequate rest and sleep to support the body's healing and recovery processes.
  \item Maintain a balanced diet with regular meals and stay well-hydrated.
  \item Monitor symptoms and track any changes in severity, frequency, or associated features.
  \item Avoid overexertion and gradually increase activity levels as tolerated.
  \item Seek medical evaluation if symptoms persist, worsen, or interfere significantly with daily function.
\end{itemize}

\section*{How Your Doctor Will Evaluate You}
\begin{itemize}
  \item Your doctor will ask about your symptoms: when they started, what triggers them, and what makes them better or worse.
  \item They will review your medical history, medications, and any recent illnesses or exposures.
  \item A physical examination will focus on the affected area and your overall health.
  \item Based on the findings, your doctor may recommend blood tests, imaging studies, or other investigations.
\end{itemize}

\section*{Treatment Options}
\begin{itemize}
  \item Treatment depends on the underlying cause identified during your evaluation.
  \item Symptom relief may include over-the-counter or prescription medications as appropriate.
  \item Lifestyle changes, rest, and home care measures can help you feel better.
  \item Your doctor will schedule follow-up visits to monitor your progress and adjust treatment as needed.
\end{itemize}

\clearpage
